% =================================================================
% SPECIFIC REGULATIONS FOR FInal Design REPORT PREPARATION
% USING LATEX
% =================================================================

\clearpage
\thispagestyle{empty}

\begin{center}
    {\fontsize{16pt}{20pt}\selectfont\bfseries
    SPECIFIC REGULATIONS FOR PREPARING\\
    FINAL DESIGN REPORTS USING \LaTeX}
\end{center}

\vspace{0.5cm}

{\fontsize{13pt}{17pt}\selectfont
\setstretch{1.3}
\justifying

% ================================================================
% I. REPORT STRUCTURE
% ================================================================

\noindent
\textbf{I.\quad REPORT STRUCTURE}

\vspace{0.15cm}

The Final Design report shall be prepared using the official
\LaTeX{} template provided by the Department. Students shall not modify
the document structure, page layout, heading styles, numbering
system, or other predefined formatting settings unless approved by
the supervisor.

The report shall be arranged in the following order:

\begin{enumerate}[label=\arabic*., leftmargin=1.2cm, itemsep=2pt]
    \item Outer cover page;
    \item Inner cover page;
    \item Final Design Assignment Form;
    \item Supervisor's Evaluation Form;
    \item Acknowledgements (if applicable);
    \item Table of Contents;
    \item List of Symbols and Abbreviations (if applicable);
    \item List of Figures;
    \item List of Tables;
    \item Abstract;
    \item Main chapters of the FInal Design;
    \item Conclusions and Recommendations;
    \item References;
    \item Appendices (if applicable).
\end{enumerate}

The preliminary pages shall not be numbered as chapters. The main
content shall be organized into consecutively numbered chapters,
beginning with Chapter 1.

\vspace{0.3cm}

% ================================================================
% II. CONTENT OF THE FInal Design
% ================================================================

\noindent
\textbf{II.\quad CONTENT OF THE FInal Design}

\vspace{0.15cm}

\begin{itemize}[leftmargin=0.8cm, itemsep=3pt]

    \item The report shall contain only information directly related
    to the assigned project topic.

    \item The content shall demonstrate the student's own work,
    understanding, analysis, design, calculation, simulation,
    implementation, and evaluation as appropriate to the project.

    \item Information obtained from books, journal papers,
    conference papers, standards, websites, technical documents,
    or other sources must be properly cited.

    \item Chapters, sections, and subsections shall be generated
    automatically using the standard \LaTeX{} commands:

\begin{verbatim}
\chapter{...}
\section{...}
\subsection{...}
\end{verbatim}

    \item Chapter, section, and subsection numbers shall not be
    manually entered.

    \item The recommended hierarchy is Chapter -- Section --
    Subsection. Excessive subdivision of the report should be avoided.

    \item Unless otherwise specified by the Faculty or supervisor,
    the main report should normally contain between 30 and 80 pages,
    excluding large appendices.

\end{itemize}

\vspace{0.3cm}

% ================================================================
% III. PAGE LAYOUT AND TEXT FORMATTING
% ================================================================

\noindent
\textbf{III.\quad PAGE LAYOUT AND TEXT FORMATTING}

\vspace{0.15cm}

The report shall be formatted on A4 paper and printed on one side
of each sheet. The page margins shall be:

\begin{center}
\begin{tabular}{ll}
    Left margin:   & 3.5 cm \\
    Right margin:  & 2.0 cm \\
    Top margin:    & 2.5 cm \\
    Bottom margin: & 2.5 cm
\end{tabular}
\end{center}

The main text shall use a font size of 13 pt with 1.3-line spacing.
A Times New Roman or equivalent Times-family typeface shall be used.

The text shall be fully justified. Students shall not use repeated
spaces, manual line breaks, or repeated blank lines to position
text.

Page layout and general formatting shall be controlled by the
official template. Students shall not manually alter margins, font
sizes, line spacing, heading spacing, or paragraph formatting for
individual pages.

\vspace{0.3cm}

% ================================================================
% IV. CHAPTERS, SECTIONS, AND SUBSECTIONS
% ================================================================

\noindent
\textbf{IV.\quad CHAPTERS, SECTIONS, AND SUBSECTIONS}

\vspace{0.15cm}

\begin{itemize}[leftmargin=0.8cm, itemsep=3pt]

    \item Chapter titles shall be generated using:

\begin{verbatim}
\chapter{Chapter Title}
\end{verbatim}

    \item Section titles shall be generated using:

\begin{verbatim}
\section{Section Title}
\end{verbatim}

    \item Subsection titles shall be generated using:

\begin{verbatim}
\subsection{Subsection Title}
\end{verbatim}

    \item Numbers such as Chapter 1, 1.1, 1.1.1, etc.\ shall be
    generated automatically by \LaTeX{}.

    \item Students shall not manually reproduce heading styles using
    commands such as \verb|\textbf|, \verb|\fontsize|,
    or \verb|\centering| instead of the appropriate structural
    command.

\end{itemize}

\vspace{0.3cm}

% ================================================================
% V. PAGE NUMBERING
% ================================================================

\noindent
\textbf{V.\quad PAGE NUMBERING}

\vspace{0.15cm}

Cover pages and official forms at the beginning of the report shall
not display page numbers.

The main body of the report shall use Arabic page numbers beginning
from Chapter 1.

Page numbers shall be generated automatically by \LaTeX{}.
Students shall not manually enter page numbers.

\vspace{0.3cm}

% ================================================================
% VI. TABLE OF CONTENTS AND LISTS
% ================================================================

\noindent
\textbf{VI.\quad TABLE OF CONTENTS AND LISTS}

\vspace{0.15cm}

The Table of Contents shall be generated automatically using:

\begin{verbatim}
\tableofcontents
\end{verbatim}

The List of Figures shall be generated using:

\begin{verbatim}
\listoffigures
\end{verbatim}

The List of Tables shall be generated using:

\begin{verbatim}
\listoftables
\end{verbatim}

Students shall not manually create or edit page numbers in the
Table of Contents, List of Figures, or List of Tables.

The document shall be compiled sufficiently to ensure that all page
numbers, references, and lists are correctly updated.

\vspace{0.3cm}

% ================================================================
% VII. FIGURES, DIAGRAMS, AND GRAPHS
% ================================================================

\noindent
\textbf{VII.\quad FIGURES, DIAGRAMS, AND GRAPHS}

\vspace{0.15cm}

\begin{itemize}[leftmargin=0.8cm, itemsep=3pt]

    \item Figures shall be clear, legible, and directly relevant to
    the report content.

    \item Figures shall be inserted using the standard
    \verb|figure| environment.

    \item Figure captions shall be placed below the figure.

    \item Figures shall be numbered automatically according to the
    numbering system defined in the template.

    \item Each figure that is referenced in the text shall be assigned
    a unique label.

\end{itemize}

Example:

\begin{verbatim}
\begin{figure}[H]
    \centering
    \includegraphics[width=0.8\textwidth]
                    {figures/system_model.pdf}
    \caption{System model}
    \label{fig:system-model}
\end{figure}
\end{verbatim}

Figures shall be referenced in the text using:

\begin{verbatim}
Figure~\ref{fig:system-model}
\end{verbatim}

Students shall not manually type figure numbers.

Figures reproduced or adapted from other sources shall include an
appropriate citation.

Whenever possible, vector graphics or high-resolution images should
be used. Low-quality screenshots should be avoided.

\vspace{0.3cm}

% ================================================================
% VIII. TABLES
% ================================================================

\noindent
\textbf{VIII.\quad TABLES}

\vspace{0.15cm}

Tables shall be prepared using appropriate \LaTeX{} environments
such as \verb|table|, \verb|tabular|, \verb|tabularx|, or
\verb|longtable|.

Table captions shall be placed above the table.

Example:

\begin{verbatim}
\begin{table}[H]
    \centering
    \caption{System parameters}
    \label{tab:system-parameters}

    \begin{tabular}{ccc}
        \hline
        Parameter & Symbol & Value \\
        \hline
        Resistance & $R$ & 0.7 $\Omega$ \\
        Inductance & $L$ & 2.72 mH \\
        \hline
    \end{tabular}
\end{table}
\end{verbatim}

Tables shall be referenced using:

\begin{verbatim}
Table~\ref{tab:system-parameters}
\end{verbatim}

Students shall not manually enter table numbers.

Tables should be created directly in \LaTeX{} rather than inserted
as screenshots whenever practicable.

\vspace{0.3cm}

% ================================================================
% IX. EQUATIONS AND MATHEMATICAL EXPRESSIONS
% ================================================================

\noindent
\textbf{IX.\quad EQUATIONS AND MATHEMATICAL EXPRESSIONS}

\vspace{0.15cm}

Mathematical expressions shall be typeset directly using \LaTeX{}.
Images or screenshots of equations shall not be used.

Short mathematical expressions appearing within a sentence shall
use inline mathematical mode.

Example:

\begin{verbatim}
$R = 0.7\,\Omega$
\end{verbatim}

Important equations shall be placed in a numbered equation
environment.

Example:

\begin{verbatim}
\begin{equation}
    J\frac{d\omega}{dt}
    = T_e - T_L - B\omega
    \label{eq:mechanical-model}
\end{equation}
\end{verbatim}

Equations shall be referenced using:

\begin{verbatim}
Equation~\eqref{eq:mechanical-model}
\end{verbatim}

Equation numbers shall not be manually entered.

Mathematical notation shall be consistent throughout the entire
report.

\vspace{0.3cm}

% ================================================================
% X. SYMBOLS, UNITS, AND ABBREVIATIONS
% ================================================================

\noindent
\textbf{X.\quad SYMBOLS, UNITS, AND ABBREVIATIONS}

\vspace{0.15cm}

\begin{itemize}[leftmargin=0.8cm, itemsep=3pt]

    \item SI units shall be used wherever applicable.

    \item Symbols and units shall be written consistently throughout
    the report.

    \item The \verb|siunitx| package should be used for physical
    quantities and units where appropriate.

\end{itemize}

Examples:

\begin{verbatim}
\SI{220}{\volt}
\SI{2.72}{\milli\henry}
\SI{450}{\radian\per\second}
\end{verbatim}

Symbols and abbreviations frequently used in the report shall be
included in the List of Symbols and Abbreviations.

Each abbreviation should normally be defined when it first appears
in the text.

\vspace{0.3cm}

% ================================================================
% XI. SOURCE CODE
% ================================================================

\noindent
\textbf{XI.\quad SOURCE CODE}

\vspace{0.15cm}

Source code that is necessary for explaining an algorithm,
controller, simulation, or implementation may be included in the
main report.

The \verb|listings| environment, or another environment defined by
the official template, shall be used.

Example:

\begin{verbatim}
\begin{lstlisting}[language=Matlab]
A = [-R/L 0;
      0  -R/L];

K = hinfsyn(P,nmeas,ncont);
\end{lstlisting}
\end{verbatim}

Long source-code listings that are not essential to the immediate
discussion should be placed in the appendices.

Source code should not normally be inserted as screenshots.

\vspace{0.3cm}

% ================================================================
% XII. CROSS-REFERENCING
% ================================================================

\noindent
\textbf{XII.\quad CROSS-REFERENCING}

\vspace{0.15cm}

Chapters, sections, figures, tables, equations, and other numbered
elements shall be cross-referenced using the standard \LaTeX{}
mechanism.

The following commands shall be used where appropriate:

\begin{verbatim}
\label{...}
\ref{...}
\eqref{...}
\end{verbatim}

For example:

\begin{verbatim}
As shown in Figure~\ref{fig:response}, ...
\end{verbatim}

Students shall not manually type a reference number such as
``Figure 3.2'' when the number can be generated automatically.

Cross-referencing ensures that all numbers remain correct when
figures, tables, equations, or sections are added or removed.

\vspace{0.3cm}

% ================================================================
% XIII. REFERENCES AND CITATIONS
% ================================================================

\noindent
\textbf{XIII.\quad REFERENCES AND CITATIONS}

\vspace{0.15cm}

All sources used in preparing the FInal Design shall be properly
cited.

References shall be managed using the bibliography system provided
with the official template. Bibliographic information shall normally
be stored in the file:

\begin{verbatim}
References.bib
\end{verbatim}

A reference may be cited in the text using:

\begin{verbatim}
\cite{reference-key}
\end{verbatim}

Students shall not manually enter reference numbers such as
``[1]'', ``[2]'', etc.

Bibliographic entries shall contain sufficient information to
identify the original source.

Titles of books, journal articles, conference papers, standards,
and other publications in foreign languages shall be retained in
their original language and shall not be translated solely for the
reference list.

All ideas, data, figures, tables, equations, algorithms, or results
obtained from another source shall be appropriately acknowledged.

\vspace{0.3cm}

% ================================================================
% XIV. APPENDICES
% ================================================================

\noindent
\textbf{XIV.\quad APPENDICES}

\vspace{0.15cm}

Supporting material that is useful but too extensive for the main
body of the report should be placed in the appendices.

Appendices may include:

\begin{itemize}[leftmargin=0.8cm, itemsep=2pt]
    \item complete source code;
    \item detailed calculation sheets;
    \item large data tables;
    \item technical drawings;
    \item additional simulation results;
    \item equipment specifications;
    \item supplementary design documents.
\end{itemize}

Appendices shall be created using the appendix environment defined
in the official template.

Example:

\begin{verbatim}
\begin{appendices}

\chapter{MATLAB SOURCE CODE}

...

\chapter{SYSTEM PARAMETERS}

...

\end{appendices}
\end{verbatim}

Appendix titles and numbering shall not be manually entered.

\vspace{0.3cm}

% ================================================================
% XV. LATEX PROJECT ORGANIZATION
% ================================================================

\noindent
\textbf{XV.\quad \LaTeX{} PROJECT ORGANIZATION}

\vspace{0.15cm}

Students shall use the official \LaTeX{} template provided by the
Faculty.

The main formatting commands contained in the document preamble
shall not be modified unless approval is obtained from the
supervisor.

A recommended project structure is:

\begin{verbatim}
FinalDesign/
|
|-- main.tex
|-- References.bib
|
|-- figures/
|   |-- figure01.pdf
|   |-- figure02.png
|   `-- ...
|
|
`-- appendix/
    `-- ...
\end{verbatim}

File and folder names should preferably contain Latin letters,
numbers, hyphens, or underscores.

Spaces, Vietnamese diacritics, and unnecessary special characters
in file names should be avoided where possible.

Recommended examples include:

\begin{verbatim}
motor_model.pdf
chapter_01.tex
simulation_result.png
\end{verbatim}

\vspace{0.3cm}

% ================================================================
% XVI. DOCUMENT COMPILATION AND FINAL CHECK
% ================================================================

\noindent
\textbf{XVI.\quad DOCUMENT COMPILATION AND FINAL CHECK}

\vspace{0.15cm}

Before submission, the complete \LaTeX{} project shall compile
successfully and produce the final PDF report.

The final document shall not contain unresolved references such as:

\begin{verbatim}
??
[?]
\end{verbatim}

Before generating the final PDF, students shall check:

\begin{itemize}[leftmargin=0.8cm, itemsep=2pt]
    \item the Table of Contents;
    \item page numbering;
    \item chapter and section numbering;
    \item figure numbering;
    \item table numbering;
    \item equation numbering;
    \item cross-references;
    \item citations;
    \item bibliography;
    \item appendices;
    \item font display;
    \item mathematical symbols;
    \item image quality.
\end{itemize}

The document may be compiled using TeX Live, MiKTeX, or
another compatible \LaTeX{} environment. Regardless of the
compilation environment, the resulting PDF shall conform to the
official Department template. \textbf{Students are not allowed to use Overleaf 
or other online \LaTeX{} editors to protect the 
knowledge and intellectual property of the project.}

\vspace{0.3cm}

% ================================================================
% XVII. SUBMISSION REQUIREMENTS
% ================================================================

\noindent
\textbf{XVII.\quad SUBMISSION REQUIREMENTS}

\vspace{0.15cm}

The printed report shall satisfy the following requirements:

\begin{itemize}[leftmargin=0.8cm, itemsep=2pt]
    \item A4 paper;
    \item single-sided printing;
    \item one bound copy, unless otherwise specified;
    \item soft cover;
    \item formatting in accordance with the official template.
\end{itemize}

The electronic submission should include:

\begin{itemize}[leftmargin=0.8cm, itemsep=2pt]
    \item the final PDF report;
    \item the complete \LaTeX{} source files;
    \item the \verb|References.bib| file;
    \item all figures and graphics;
    \item all additional files required to compile the document.
\end{itemize}

The submitted source project shall be complete and capable of being
recompiled to produce a PDF corresponding to the submitted report.

\vspace{0.3cm}

% ================================================================
% XVIII. ELIGIBILITY FOR FInal Design DEFENSE
% ================================================================

\noindent
\textbf{XVIII.\quad ELIGIBILITY FOR FINAL DESIGN DEFENSE}

\vspace{0.15cm}

A student shall be eligible to defend the FInal Design when the
following requirements have been satisfied:

\begin{enumerate}[label=\arabic*., leftmargin=1.2cm, itemsep=2pt]

    \item The FInal Design has been completed within the specified
    period.

    \item The report satisfies the academic and technical requirements
    of the assigned project.

    \item The report complies with the official \LaTeX{} template
    issued by the Faculty.

    \item The complete \LaTeX{} source compiles successfully without
    errors affecting the document content or presentation.

    \item The final PDF and required source files have been completed
    and submitted.

    \item All external sources have been appropriately cited and the
    reference list has been properly generated.

    \item The supervisor has approved the student to proceed to the
    FInal Design defense.

\end{enumerate}

\vspace{0.4cm}

\noindent
\textbf{Note:}
Students are responsible for the academic and technical content of
the report. The official \LaTeX{} template is responsible for the
standardized document formatting. Students should therefore focus
on entering and organizing the content correctly rather than
manually modifying the predefined formatting parameters. 

\noindent
\textbf{Note:}
This section is provided for guidance only and is included in the main
document using the command \verb|% =================================================================
% SPECIFIC REGULATIONS FOR FInal Design REPORT PREPARATION
% USING LATEX
% =================================================================

\clearpage
\thispagestyle{empty}

\begin{center}
    {\fontsize{16pt}{20pt}\selectfont\bfseries
    SPECIFIC REGULATIONS FOR PREPARING\\
    FINAL DESIGN REPORTS USING \LaTeX}
\end{center}

\vspace{0.5cm}

{\fontsize{13pt}{17pt}\selectfont
\setstretch{1.3}
\justifying

% ================================================================
% I. REPORT STRUCTURE
% ================================================================

\noindent
\textbf{I.\quad REPORT STRUCTURE}

\vspace{0.15cm}

The Final Design report shall be prepared using the official
\LaTeX{} template provided by the Department. Students shall not modify
the document structure, page layout, heading styles, numbering
system, or other predefined formatting settings unless approved by
the supervisor.

The report shall be arranged in the following order:

\begin{enumerate}[label=\arabic*., leftmargin=1.2cm, itemsep=2pt]
    \item Outer cover page;
    \item Inner cover page;
    \item Final Design Assignment Form;
    \item Supervisor's Evaluation Form;
    \item Acknowledgements (if applicable);
    \item Table of Contents;
    \item List of Symbols and Abbreviations (if applicable);
    \item List of Figures;
    \item List of Tables;
    \item Abstract;
    \item Main chapters of the FInal Design;
    \item Conclusions and Recommendations;
    \item References;
    \item Appendices (if applicable).
\end{enumerate}

The preliminary pages shall not be numbered as chapters. The main
content shall be organized into consecutively numbered chapters,
beginning with Chapter 1.

\vspace{0.3cm}

% ================================================================
% II. CONTENT OF THE FInal Design
% ================================================================

\noindent
\textbf{II.\quad CONTENT OF THE FInal Design}

\vspace{0.15cm}

\begin{itemize}[leftmargin=0.8cm, itemsep=3pt]

    \item The report shall contain only information directly related
    to the assigned project topic.

    \item The content shall demonstrate the student's own work,
    understanding, analysis, design, calculation, simulation,
    implementation, and evaluation as appropriate to the project.

    \item Information obtained from books, journal papers,
    conference papers, standards, websites, technical documents,
    or other sources must be properly cited.

    \item Chapters, sections, and subsections shall be generated
    automatically using the standard \LaTeX{} commands:

\begin{verbatim}
\chapter{...}
\section{...}
\subsection{...}
\end{verbatim}

    \item Chapter, section, and subsection numbers shall not be
    manually entered.

    \item The recommended hierarchy is Chapter -- Section --
    Subsection. Excessive subdivision of the report should be avoided.

    \item Unless otherwise specified by the Faculty or supervisor,
    the main report should normally contain between 30 and 80 pages,
    excluding large appendices.

\end{itemize}

\vspace{0.3cm}

% ================================================================
% III. PAGE LAYOUT AND TEXT FORMATTING
% ================================================================

\noindent
\textbf{III.\quad PAGE LAYOUT AND TEXT FORMATTING}

\vspace{0.15cm}

The report shall be formatted on A4 paper and printed on one side
of each sheet. The page margins shall be:

\begin{center}
\begin{tabular}{ll}
    Left margin:   & 3.5 cm \\
    Right margin:  & 2.0 cm \\
    Top margin:    & 2.5 cm \\
    Bottom margin: & 2.5 cm
\end{tabular}
\end{center}

The main text shall use a font size of 13 pt with 1.3-line spacing.
A Times New Roman or equivalent Times-family typeface shall be used.

The text shall be fully justified. Students shall not use repeated
spaces, manual line breaks, or repeated blank lines to position
text.

Page layout and general formatting shall be controlled by the
official template. Students shall not manually alter margins, font
sizes, line spacing, heading spacing, or paragraph formatting for
individual pages.

\vspace{0.3cm}

% ================================================================
% IV. CHAPTERS, SECTIONS, AND SUBSECTIONS
% ================================================================

\noindent
\textbf{IV.\quad CHAPTERS, SECTIONS, AND SUBSECTIONS}

\vspace{0.15cm}

\begin{itemize}[leftmargin=0.8cm, itemsep=3pt]

    \item Chapter titles shall be generated using:

\begin{verbatim}
\chapter{Chapter Title}
\end{verbatim}

    \item Section titles shall be generated using:

\begin{verbatim}
\section{Section Title}
\end{verbatim}

    \item Subsection titles shall be generated using:

\begin{verbatim}
\subsection{Subsection Title}
\end{verbatim}

    \item Numbers such as Chapter 1, 1.1, 1.1.1, etc.\ shall be
    generated automatically by \LaTeX{}.

    \item Students shall not manually reproduce heading styles using
    commands such as \verb|\textbf|, \verb|\fontsize|,
    or \verb|\centering| instead of the appropriate structural
    command.

\end{itemize}

\vspace{0.3cm}

% ================================================================
% V. PAGE NUMBERING
% ================================================================

\noindent
\textbf{V.\quad PAGE NUMBERING}

\vspace{0.15cm}

Cover pages and official forms at the beginning of the report shall
not display page numbers.

The main body of the report shall use Arabic page numbers beginning
from Chapter 1.

Page numbers shall be generated automatically by \LaTeX{}.
Students shall not manually enter page numbers.

\vspace{0.3cm}

% ================================================================
% VI. TABLE OF CONTENTS AND LISTS
% ================================================================

\noindent
\textbf{VI.\quad TABLE OF CONTENTS AND LISTS}

\vspace{0.15cm}

The Table of Contents shall be generated automatically using:

\begin{verbatim}
\tableofcontents
\end{verbatim}

The List of Figures shall be generated using:

\begin{verbatim}
\listoffigures
\end{verbatim}

The List of Tables shall be generated using:

\begin{verbatim}
\listoftables
\end{verbatim}

Students shall not manually create or edit page numbers in the
Table of Contents, List of Figures, or List of Tables.

The document shall be compiled sufficiently to ensure that all page
numbers, references, and lists are correctly updated.

\vspace{0.3cm}

% ================================================================
% VII. FIGURES, DIAGRAMS, AND GRAPHS
% ================================================================

\noindent
\textbf{VII.\quad FIGURES, DIAGRAMS, AND GRAPHS}

\vspace{0.15cm}

\begin{itemize}[leftmargin=0.8cm, itemsep=3pt]

    \item Figures shall be clear, legible, and directly relevant to
    the report content.

    \item Figures shall be inserted using the standard
    \verb|figure| environment.

    \item Figure captions shall be placed below the figure.

    \item Figures shall be numbered automatically according to the
    numbering system defined in the template.

    \item Each figure that is referenced in the text shall be assigned
    a unique label.

\end{itemize}

Example:

\begin{verbatim}
\begin{figure}[H]
    \centering
    \includegraphics[width=0.8\textwidth]
                    {figures/system_model.pdf}
    \caption{System model}
    \label{fig:system-model}
\end{figure}
\end{verbatim}

Figures shall be referenced in the text using:

\begin{verbatim}
Figure~\ref{fig:system-model}
\end{verbatim}

Students shall not manually type figure numbers.

Figures reproduced or adapted from other sources shall include an
appropriate citation.

Whenever possible, vector graphics or high-resolution images should
be used. Low-quality screenshots should be avoided.

\vspace{0.3cm}

% ================================================================
% VIII. TABLES
% ================================================================

\noindent
\textbf{VIII.\quad TABLES}

\vspace{0.15cm}

Tables shall be prepared using appropriate \LaTeX{} environments
such as \verb|table|, \verb|tabular|, \verb|tabularx|, or
\verb|longtable|.

Table captions shall be placed above the table.

Example:

\begin{verbatim}
\begin{table}[H]
    \centering
    \caption{System parameters}
    \label{tab:system-parameters}

    \begin{tabular}{ccc}
        \hline
        Parameter & Symbol & Value \\
        \hline
        Resistance & $R$ & 0.7 $\Omega$ \\
        Inductance & $L$ & 2.72 mH \\
        \hline
    \end{tabular}
\end{table}
\end{verbatim}

Tables shall be referenced using:

\begin{verbatim}
Table~\ref{tab:system-parameters}
\end{verbatim}

Students shall not manually enter table numbers.

Tables should be created directly in \LaTeX{} rather than inserted
as screenshots whenever practicable.

\vspace{0.3cm}

% ================================================================
% IX. EQUATIONS AND MATHEMATICAL EXPRESSIONS
% ================================================================

\noindent
\textbf{IX.\quad EQUATIONS AND MATHEMATICAL EXPRESSIONS}

\vspace{0.15cm}

Mathematical expressions shall be typeset directly using \LaTeX{}.
Images or screenshots of equations shall not be used.

Short mathematical expressions appearing within a sentence shall
use inline mathematical mode.

Example:

\begin{verbatim}
$R = 0.7\,\Omega$
\end{verbatim}

Important equations shall be placed in a numbered equation
environment.

Example:

\begin{verbatim}
\begin{equation}
    J\frac{d\omega}{dt}
    = T_e - T_L - B\omega
    \label{eq:mechanical-model}
\end{equation}
\end{verbatim}

Equations shall be referenced using:

\begin{verbatim}
Equation~\eqref{eq:mechanical-model}
\end{verbatim}

Equation numbers shall not be manually entered.

Mathematical notation shall be consistent throughout the entire
report.

\vspace{0.3cm}

% ================================================================
% X. SYMBOLS, UNITS, AND ABBREVIATIONS
% ================================================================

\noindent
\textbf{X.\quad SYMBOLS, UNITS, AND ABBREVIATIONS}

\vspace{0.15cm}

\begin{itemize}[leftmargin=0.8cm, itemsep=3pt]

    \item SI units shall be used wherever applicable.

    \item Symbols and units shall be written consistently throughout
    the report.

    \item The \verb|siunitx| package should be used for physical
    quantities and units where appropriate.

\end{itemize}

Examples:

\begin{verbatim}
\SI{220}{\volt}
\SI{2.72}{\milli\henry}
\SI{450}{\radian\per\second}
\end{verbatim}

Symbols and abbreviations frequently used in the report shall be
included in the List of Symbols and Abbreviations.

Each abbreviation should normally be defined when it first appears
in the text.

\vspace{0.3cm}

% ================================================================
% XI. SOURCE CODE
% ================================================================

\noindent
\textbf{XI.\quad SOURCE CODE}

\vspace{0.15cm}

Source code that is necessary for explaining an algorithm,
controller, simulation, or implementation may be included in the
main report.

The \verb|listings| environment, or another environment defined by
the official template, shall be used.

Example:

\begin{verbatim}
\begin{lstlisting}[language=Matlab]
A = [-R/L 0;
      0  -R/L];

K = hinfsyn(P,nmeas,ncont);
\end{lstlisting}
\end{verbatim}

Long source-code listings that are not essential to the immediate
discussion should be placed in the appendices.

Source code should not normally be inserted as screenshots.

\vspace{0.3cm}

% ================================================================
% XII. CROSS-REFERENCING
% ================================================================

\noindent
\textbf{XII.\quad CROSS-REFERENCING}

\vspace{0.15cm}

Chapters, sections, figures, tables, equations, and other numbered
elements shall be cross-referenced using the standard \LaTeX{}
mechanism.

The following commands shall be used where appropriate:

\begin{verbatim}
\label{...}
\ref{...}
\eqref{...}
\end{verbatim}

For example:

\begin{verbatim}
As shown in Figure~\ref{fig:response}, ...
\end{verbatim}

Students shall not manually type a reference number such as
``Figure 3.2'' when the number can be generated automatically.

Cross-referencing ensures that all numbers remain correct when
figures, tables, equations, or sections are added or removed.

\vspace{0.3cm}

% ================================================================
% XIII. REFERENCES AND CITATIONS
% ================================================================

\noindent
\textbf{XIII.\quad REFERENCES AND CITATIONS}

\vspace{0.15cm}

All sources used in preparing the FInal Design shall be properly
cited.

References shall be managed using the bibliography system provided
with the official template. Bibliographic information shall normally
be stored in the file:

\begin{verbatim}
References.bib
\end{verbatim}

A reference may be cited in the text using:

\begin{verbatim}
\cite{reference-key}
\end{verbatim}

Students shall not manually enter reference numbers such as
``[1]'', ``[2]'', etc.

Bibliographic entries shall contain sufficient information to
identify the original source.

Titles of books, journal articles, conference papers, standards,
and other publications in foreign languages shall be retained in
their original language and shall not be translated solely for the
reference list.

All ideas, data, figures, tables, equations, algorithms, or results
obtained from another source shall be appropriately acknowledged.

\vspace{0.3cm}

% ================================================================
% XIV. APPENDICES
% ================================================================

\noindent
\textbf{XIV.\quad APPENDICES}

\vspace{0.15cm}

Supporting material that is useful but too extensive for the main
body of the report should be placed in the appendices.

Appendices may include:

\begin{itemize}[leftmargin=0.8cm, itemsep=2pt]
    \item complete source code;
    \item detailed calculation sheets;
    \item large data tables;
    \item technical drawings;
    \item additional simulation results;
    \item equipment specifications;
    \item supplementary design documents.
\end{itemize}

Appendices shall be created using the appendix environment defined
in the official template.

Example:

\begin{verbatim}
\begin{appendices}

\chapter{MATLAB SOURCE CODE}

...

\chapter{SYSTEM PARAMETERS}

...

\end{appendices}
\end{verbatim}

Appendix titles and numbering shall not be manually entered.

\vspace{0.3cm}

% ================================================================
% XV. LATEX PROJECT ORGANIZATION
% ================================================================

\noindent
\textbf{XV.\quad \LaTeX{} PROJECT ORGANIZATION}

\vspace{0.15cm}

Students shall use the official \LaTeX{} template provided by the
Faculty.

The main formatting commands contained in the document preamble
shall not be modified unless approval is obtained from the
supervisor.

A recommended project structure is:

\begin{verbatim}
FinalDesign/
|
|-- main.tex
|-- References.bib
|
|-- figures/
|   |-- figure01.pdf
|   |-- figure02.png
|   `-- ...
|
|
`-- appendix/
    `-- ...
\end{verbatim}

File and folder names should preferably contain Latin letters,
numbers, hyphens, or underscores.

Spaces, Vietnamese diacritics, and unnecessary special characters
in file names should be avoided where possible.

Recommended examples include:

\begin{verbatim}
motor_model.pdf
chapter_01.tex
simulation_result.png
\end{verbatim}

\vspace{0.3cm}

% ================================================================
% XVI. DOCUMENT COMPILATION AND FINAL CHECK
% ================================================================

\noindent
\textbf{XVI.\quad DOCUMENT COMPILATION AND FINAL CHECK}

\vspace{0.15cm}

Before submission, the complete \LaTeX{} project shall compile
successfully and produce the final PDF report.

The final document shall not contain unresolved references such as:

\begin{verbatim}
??
[?]
\end{verbatim}

Before generating the final PDF, students shall check:

\begin{itemize}[leftmargin=0.8cm, itemsep=2pt]
    \item the Table of Contents;
    \item page numbering;
    \item chapter and section numbering;
    \item figure numbering;
    \item table numbering;
    \item equation numbering;
    \item cross-references;
    \item citations;
    \item bibliography;
    \item appendices;
    \item font display;
    \item mathematical symbols;
    \item image quality.
\end{itemize}

The document may be compiled using TeX Live, MiKTeX, or
another compatible \LaTeX{} environment. Regardless of the
compilation environment, the resulting PDF shall conform to the
official Department template. \textbf{Students are not allowed to use Overleaf 
or other online \LaTeX{} editors to protect the 
knowledge and intellectual property of the project.}

\vspace{0.3cm}

% ================================================================
% XVII. SUBMISSION REQUIREMENTS
% ================================================================

\noindent
\textbf{XVII.\quad SUBMISSION REQUIREMENTS}

\vspace{0.15cm}

The printed report shall satisfy the following requirements:

\begin{itemize}[leftmargin=0.8cm, itemsep=2pt]
    \item A4 paper;
    \item single-sided printing;
    \item one bound copy, unless otherwise specified;
    \item soft cover;
    \item formatting in accordance with the official template.
\end{itemize}

The electronic submission should include:

\begin{itemize}[leftmargin=0.8cm, itemsep=2pt]
    \item the final PDF report;
    \item the complete \LaTeX{} source files;
    \item the \verb|References.bib| file;
    \item all figures and graphics;
    \item all additional files required to compile the document.
\end{itemize}

The submitted source project shall be complete and capable of being
recompiled to produce a PDF corresponding to the submitted report.

\vspace{0.3cm}

% ================================================================
% XVIII. ELIGIBILITY FOR FInal Design DEFENSE
% ================================================================

\noindent
\textbf{XVIII.\quad ELIGIBILITY FOR FINAL DESIGN DEFENSE}

\vspace{0.15cm}

A student shall be eligible to defend the FInal Design when the
following requirements have been satisfied:

\begin{enumerate}[label=\arabic*., leftmargin=1.2cm, itemsep=2pt]

    \item The FInal Design has been completed within the specified
    period.

    \item The report satisfies the academic and technical requirements
    of the assigned project.

    \item The report complies with the official \LaTeX{} template
    issued by the Faculty.

    \item The complete \LaTeX{} source compiles successfully without
    errors affecting the document content or presentation.

    \item The final PDF and required source files have been completed
    and submitted.

    \item All external sources have been appropriately cited and the
    reference list has been properly generated.

    \item The supervisor has approved the student to proceed to the
    FInal Design defense.

\end{enumerate}

\vspace{0.4cm}

\noindent
\textbf{Note:}
Students are responsible for the academic and technical content of
the report. The official \LaTeX{} template is responsible for the
standardized document formatting. Students should therefore focus
on entering and organizing the content correctly rather than
manually modifying the predefined formatting parameters. 

\noindent
\textbf{Note:}
This section is provided for guidance only and is included in the main
document using the command \verb|% =================================================================
% SPECIFIC REGULATIONS FOR FInal Design REPORT PREPARATION
% USING LATEX
% =================================================================

\clearpage
\thispagestyle{empty}

\begin{center}
    {\fontsize{16pt}{20pt}\selectfont\bfseries
    SPECIFIC REGULATIONS FOR PREPARING\\
    FINAL DESIGN REPORTS USING \LaTeX}
\end{center}

\vspace{0.5cm}

{\fontsize{13pt}{17pt}\selectfont
\setstretch{1.3}
\justifying

% ================================================================
% I. REPORT STRUCTURE
% ================================================================

\noindent
\textbf{I.\quad REPORT STRUCTURE}

\vspace{0.15cm}

The Final Design report shall be prepared using the official
\LaTeX{} template provided by the Department. Students shall not modify
the document structure, page layout, heading styles, numbering
system, or other predefined formatting settings unless approved by
the supervisor.

The report shall be arranged in the following order:

\begin{enumerate}[label=\arabic*., leftmargin=1.2cm, itemsep=2pt]
    \item Outer cover page;
    \item Inner cover page;
    \item Final Design Assignment Form;
    \item Supervisor's Evaluation Form;
    \item Acknowledgements (if applicable);
    \item Table of Contents;
    \item List of Symbols and Abbreviations (if applicable);
    \item List of Figures;
    \item List of Tables;
    \item Abstract;
    \item Main chapters of the FInal Design;
    \item Conclusions and Recommendations;
    \item References;
    \item Appendices (if applicable).
\end{enumerate}

The preliminary pages shall not be numbered as chapters. The main
content shall be organized into consecutively numbered chapters,
beginning with Chapter 1.

\vspace{0.3cm}

% ================================================================
% II. CONTENT OF THE FInal Design
% ================================================================

\noindent
\textbf{II.\quad CONTENT OF THE FInal Design}

\vspace{0.15cm}

\begin{itemize}[leftmargin=0.8cm, itemsep=3pt]

    \item The report shall contain only information directly related
    to the assigned project topic.

    \item The content shall demonstrate the student's own work,
    understanding, analysis, design, calculation, simulation,
    implementation, and evaluation as appropriate to the project.

    \item Information obtained from books, journal papers,
    conference papers, standards, websites, technical documents,
    or other sources must be properly cited.

    \item Chapters, sections, and subsections shall be generated
    automatically using the standard \LaTeX{} commands:

\begin{verbatim}
\chapter{...}
\section{...}
\subsection{...}
\end{verbatim}

    \item Chapter, section, and subsection numbers shall not be
    manually entered.

    \item The recommended hierarchy is Chapter -- Section --
    Subsection. Excessive subdivision of the report should be avoided.

    \item Unless otherwise specified by the Faculty or supervisor,
    the main report should normally contain between 30 and 80 pages,
    excluding large appendices.

\end{itemize}

\vspace{0.3cm}

% ================================================================
% III. PAGE LAYOUT AND TEXT FORMATTING
% ================================================================

\noindent
\textbf{III.\quad PAGE LAYOUT AND TEXT FORMATTING}

\vspace{0.15cm}

The report shall be formatted on A4 paper and printed on one side
of each sheet. The page margins shall be:

\begin{center}
\begin{tabular}{ll}
    Left margin:   & 3.5 cm \\
    Right margin:  & 2.0 cm \\
    Top margin:    & 2.5 cm \\
    Bottom margin: & 2.5 cm
\end{tabular}
\end{center}

The main text shall use a font size of 13 pt with 1.3-line spacing.
A Times New Roman or equivalent Times-family typeface shall be used.

The text shall be fully justified. Students shall not use repeated
spaces, manual line breaks, or repeated blank lines to position
text.

Page layout and general formatting shall be controlled by the
official template. Students shall not manually alter margins, font
sizes, line spacing, heading spacing, or paragraph formatting for
individual pages.

\vspace{0.3cm}

% ================================================================
% IV. CHAPTERS, SECTIONS, AND SUBSECTIONS
% ================================================================

\noindent
\textbf{IV.\quad CHAPTERS, SECTIONS, AND SUBSECTIONS}

\vspace{0.15cm}

\begin{itemize}[leftmargin=0.8cm, itemsep=3pt]

    \item Chapter titles shall be generated using:

\begin{verbatim}
\chapter{Chapter Title}
\end{verbatim}

    \item Section titles shall be generated using:

\begin{verbatim}
\section{Section Title}
\end{verbatim}

    \item Subsection titles shall be generated using:

\begin{verbatim}
\subsection{Subsection Title}
\end{verbatim}

    \item Numbers such as Chapter 1, 1.1, 1.1.1, etc.\ shall be
    generated automatically by \LaTeX{}.

    \item Students shall not manually reproduce heading styles using
    commands such as \verb|\textbf|, \verb|\fontsize|,
    or \verb|\centering| instead of the appropriate structural
    command.

\end{itemize}

\vspace{0.3cm}

% ================================================================
% V. PAGE NUMBERING
% ================================================================

\noindent
\textbf{V.\quad PAGE NUMBERING}

\vspace{0.15cm}

Cover pages and official forms at the beginning of the report shall
not display page numbers.

The main body of the report shall use Arabic page numbers beginning
from Chapter 1.

Page numbers shall be generated automatically by \LaTeX{}.
Students shall not manually enter page numbers.

\vspace{0.3cm}

% ================================================================
% VI. TABLE OF CONTENTS AND LISTS
% ================================================================

\noindent
\textbf{VI.\quad TABLE OF CONTENTS AND LISTS}

\vspace{0.15cm}

The Table of Contents shall be generated automatically using:

\begin{verbatim}
\tableofcontents
\end{verbatim}

The List of Figures shall be generated using:

\begin{verbatim}
\listoffigures
\end{verbatim}

The List of Tables shall be generated using:

\begin{verbatim}
\listoftables
\end{verbatim}

Students shall not manually create or edit page numbers in the
Table of Contents, List of Figures, or List of Tables.

The document shall be compiled sufficiently to ensure that all page
numbers, references, and lists are correctly updated.

\vspace{0.3cm}

% ================================================================
% VII. FIGURES, DIAGRAMS, AND GRAPHS
% ================================================================

\noindent
\textbf{VII.\quad FIGURES, DIAGRAMS, AND GRAPHS}

\vspace{0.15cm}

\begin{itemize}[leftmargin=0.8cm, itemsep=3pt]

    \item Figures shall be clear, legible, and directly relevant to
    the report content.

    \item Figures shall be inserted using the standard
    \verb|figure| environment.

    \item Figure captions shall be placed below the figure.

    \item Figures shall be numbered automatically according to the
    numbering system defined in the template.

    \item Each figure that is referenced in the text shall be assigned
    a unique label.

\end{itemize}

Example:

\begin{verbatim}
\begin{figure}[H]
    \centering
    \includegraphics[width=0.8\textwidth]
                    {figures/system_model.pdf}
    \caption{System model}
    \label{fig:system-model}
\end{figure}
\end{verbatim}

Figures shall be referenced in the text using:

\begin{verbatim}
Figure~\ref{fig:system-model}
\end{verbatim}

Students shall not manually type figure numbers.

Figures reproduced or adapted from other sources shall include an
appropriate citation.

Whenever possible, vector graphics or high-resolution images should
be used. Low-quality screenshots should be avoided.

\vspace{0.3cm}

% ================================================================
% VIII. TABLES
% ================================================================

\noindent
\textbf{VIII.\quad TABLES}

\vspace{0.15cm}

Tables shall be prepared using appropriate \LaTeX{} environments
such as \verb|table|, \verb|tabular|, \verb|tabularx|, or
\verb|longtable|.

Table captions shall be placed above the table.

Example:

\begin{verbatim}
\begin{table}[H]
    \centering
    \caption{System parameters}
    \label{tab:system-parameters}

    \begin{tabular}{ccc}
        \hline
        Parameter & Symbol & Value \\
        \hline
        Resistance & $R$ & 0.7 $\Omega$ \\
        Inductance & $L$ & 2.72 mH \\
        \hline
    \end{tabular}
\end{table}
\end{verbatim}

Tables shall be referenced using:

\begin{verbatim}
Table~\ref{tab:system-parameters}
\end{verbatim}

Students shall not manually enter table numbers.

Tables should be created directly in \LaTeX{} rather than inserted
as screenshots whenever practicable.

\vspace{0.3cm}

% ================================================================
% IX. EQUATIONS AND MATHEMATICAL EXPRESSIONS
% ================================================================

\noindent
\textbf{IX.\quad EQUATIONS AND MATHEMATICAL EXPRESSIONS}

\vspace{0.15cm}

Mathematical expressions shall be typeset directly using \LaTeX{}.
Images or screenshots of equations shall not be used.

Short mathematical expressions appearing within a sentence shall
use inline mathematical mode.

Example:

\begin{verbatim}
$R = 0.7\,\Omega$
\end{verbatim}

Important equations shall be placed in a numbered equation
environment.

Example:

\begin{verbatim}
\begin{equation}
    J\frac{d\omega}{dt}
    = T_e - T_L - B\omega
    \label{eq:mechanical-model}
\end{equation}
\end{verbatim}

Equations shall be referenced using:

\begin{verbatim}
Equation~\eqref{eq:mechanical-model}
\end{verbatim}

Equation numbers shall not be manually entered.

Mathematical notation shall be consistent throughout the entire
report.

\vspace{0.3cm}

% ================================================================
% X. SYMBOLS, UNITS, AND ABBREVIATIONS
% ================================================================

\noindent
\textbf{X.\quad SYMBOLS, UNITS, AND ABBREVIATIONS}

\vspace{0.15cm}

\begin{itemize}[leftmargin=0.8cm, itemsep=3pt]

    \item SI units shall be used wherever applicable.

    \item Symbols and units shall be written consistently throughout
    the report.

    \item The \verb|siunitx| package should be used for physical
    quantities and units where appropriate.

\end{itemize}

Examples:

\begin{verbatim}
\SI{220}{\volt}
\SI{2.72}{\milli\henry}
\SI{450}{\radian\per\second}
\end{verbatim}

Symbols and abbreviations frequently used in the report shall be
included in the List of Symbols and Abbreviations.

Each abbreviation should normally be defined when it first appears
in the text.

\vspace{0.3cm}

% ================================================================
% XI. SOURCE CODE
% ================================================================

\noindent
\textbf{XI.\quad SOURCE CODE}

\vspace{0.15cm}

Source code that is necessary for explaining an algorithm,
controller, simulation, or implementation may be included in the
main report.

The \verb|listings| environment, or another environment defined by
the official template, shall be used.

Example:

\begin{verbatim}
\begin{lstlisting}[language=Matlab]
A = [-R/L 0;
      0  -R/L];

K = hinfsyn(P,nmeas,ncont);
\end{lstlisting}
\end{verbatim}

Long source-code listings that are not essential to the immediate
discussion should be placed in the appendices.

Source code should not normally be inserted as screenshots.

\vspace{0.3cm}

% ================================================================
% XII. CROSS-REFERENCING
% ================================================================

\noindent
\textbf{XII.\quad CROSS-REFERENCING}

\vspace{0.15cm}

Chapters, sections, figures, tables, equations, and other numbered
elements shall be cross-referenced using the standard \LaTeX{}
mechanism.

The following commands shall be used where appropriate:

\begin{verbatim}
\label{...}
\ref{...}
\eqref{...}
\end{verbatim}

For example:

\begin{verbatim}
As shown in Figure~\ref{fig:response}, ...
\end{verbatim}

Students shall not manually type a reference number such as
``Figure 3.2'' when the number can be generated automatically.

Cross-referencing ensures that all numbers remain correct when
figures, tables, equations, or sections are added or removed.

\vspace{0.3cm}

% ================================================================
% XIII. REFERENCES AND CITATIONS
% ================================================================

\noindent
\textbf{XIII.\quad REFERENCES AND CITATIONS}

\vspace{0.15cm}

All sources used in preparing the FInal Design shall be properly
cited.

References shall be managed using the bibliography system provided
with the official template. Bibliographic information shall normally
be stored in the file:

\begin{verbatim}
References.bib
\end{verbatim}

A reference may be cited in the text using:

\begin{verbatim}
\cite{reference-key}
\end{verbatim}

Students shall not manually enter reference numbers such as
``[1]'', ``[2]'', etc.

Bibliographic entries shall contain sufficient information to
identify the original source.

Titles of books, journal articles, conference papers, standards,
and other publications in foreign languages shall be retained in
their original language and shall not be translated solely for the
reference list.

All ideas, data, figures, tables, equations, algorithms, or results
obtained from another source shall be appropriately acknowledged.

\vspace{0.3cm}

% ================================================================
% XIV. APPENDICES
% ================================================================

\noindent
\textbf{XIV.\quad APPENDICES}

\vspace{0.15cm}

Supporting material that is useful but too extensive for the main
body of the report should be placed in the appendices.

Appendices may include:

\begin{itemize}[leftmargin=0.8cm, itemsep=2pt]
    \item complete source code;
    \item detailed calculation sheets;
    \item large data tables;
    \item technical drawings;
    \item additional simulation results;
    \item equipment specifications;
    \item supplementary design documents.
\end{itemize}

Appendices shall be created using the appendix environment defined
in the official template.

Example:

\begin{verbatim}
\begin{appendices}

\chapter{MATLAB SOURCE CODE}

...

\chapter{SYSTEM PARAMETERS}

...

\end{appendices}
\end{verbatim}

Appendix titles and numbering shall not be manually entered.

\vspace{0.3cm}

% ================================================================
% XV. LATEX PROJECT ORGANIZATION
% ================================================================

\noindent
\textbf{XV.\quad \LaTeX{} PROJECT ORGANIZATION}

\vspace{0.15cm}

Students shall use the official \LaTeX{} template provided by the
Faculty.

The main formatting commands contained in the document preamble
shall not be modified unless approval is obtained from the
supervisor.

A recommended project structure is:

\begin{verbatim}
FinalDesign/
|
|-- main.tex
|-- References.bib
|
|-- figures/
|   |-- figure01.pdf
|   |-- figure02.png
|   `-- ...
|
|
`-- appendix/
    `-- ...
\end{verbatim}

File and folder names should preferably contain Latin letters,
numbers, hyphens, or underscores.

Spaces, Vietnamese diacritics, and unnecessary special characters
in file names should be avoided where possible.

Recommended examples include:

\begin{verbatim}
motor_model.pdf
chapter_01.tex
simulation_result.png
\end{verbatim}

\vspace{0.3cm}

% ================================================================
% XVI. DOCUMENT COMPILATION AND FINAL CHECK
% ================================================================

\noindent
\textbf{XVI.\quad DOCUMENT COMPILATION AND FINAL CHECK}

\vspace{0.15cm}

Before submission, the complete \LaTeX{} project shall compile
successfully and produce the final PDF report.

The final document shall not contain unresolved references such as:

\begin{verbatim}
??
[?]
\end{verbatim}

Before generating the final PDF, students shall check:

\begin{itemize}[leftmargin=0.8cm, itemsep=2pt]
    \item the Table of Contents;
    \item page numbering;
    \item chapter and section numbering;
    \item figure numbering;
    \item table numbering;
    \item equation numbering;
    \item cross-references;
    \item citations;
    \item bibliography;
    \item appendices;
    \item font display;
    \item mathematical symbols;
    \item image quality.
\end{itemize}

The document may be compiled using TeX Live, MiKTeX, or
another compatible \LaTeX{} environment. Regardless of the
compilation environment, the resulting PDF shall conform to the
official Department template. \textbf{Students are not allowed to use Overleaf 
or other online \LaTeX{} editors to protect the 
knowledge and intellectual property of the project.}

\vspace{0.3cm}

% ================================================================
% XVII. SUBMISSION REQUIREMENTS
% ================================================================

\noindent
\textbf{XVII.\quad SUBMISSION REQUIREMENTS}

\vspace{0.15cm}

The printed report shall satisfy the following requirements:

\begin{itemize}[leftmargin=0.8cm, itemsep=2pt]
    \item A4 paper;
    \item single-sided printing;
    \item one bound copy, unless otherwise specified;
    \item soft cover;
    \item formatting in accordance with the official template.
\end{itemize}

The electronic submission should include:

\begin{itemize}[leftmargin=0.8cm, itemsep=2pt]
    \item the final PDF report;
    \item the complete \LaTeX{} source files;
    \item the \verb|References.bib| file;
    \item all figures and graphics;
    \item all additional files required to compile the document.
\end{itemize}

The submitted source project shall be complete and capable of being
recompiled to produce a PDF corresponding to the submitted report.

\vspace{0.3cm}

% ================================================================
% XVIII. ELIGIBILITY FOR FInal Design DEFENSE
% ================================================================

\noindent
\textbf{XVIII.\quad ELIGIBILITY FOR FINAL DESIGN DEFENSE}

\vspace{0.15cm}

A student shall be eligible to defend the FInal Design when the
following requirements have been satisfied:

\begin{enumerate}[label=\arabic*., leftmargin=1.2cm, itemsep=2pt]

    \item The FInal Design has been completed within the specified
    period.

    \item The report satisfies the academic and technical requirements
    of the assigned project.

    \item The report complies with the official \LaTeX{} template
    issued by the Faculty.

    \item The complete \LaTeX{} source compiles successfully without
    errors affecting the document content or presentation.

    \item The final PDF and required source files have been completed
    and submitted.

    \item All external sources have been appropriately cited and the
    reference list has been properly generated.

    \item The supervisor has approved the student to proceed to the
    FInal Design defense.

\end{enumerate}

\vspace{0.4cm}

\noindent
\textbf{Note:}
Students are responsible for the academic and technical content of
the report. The official \LaTeX{} template is responsible for the
standardized document formatting. Students should therefore focus
on entering and organizing the content correctly rather than
manually modifying the predefined formatting parameters. 

\noindent
\textbf{Note:}
This section is provided for guidance only and is included in the main
document using the command \verb|% =================================================================
% SPECIFIC REGULATIONS FOR FInal Design REPORT PREPARATION
% USING LATEX
% =================================================================

\clearpage
\thispagestyle{empty}

\begin{center}
    {\fontsize{16pt}{20pt}\selectfont\bfseries
    SPECIFIC REGULATIONS FOR PREPARING\\
    FINAL DESIGN REPORTS USING \LaTeX}
\end{center}

\vspace{0.5cm}

{\fontsize{13pt}{17pt}\selectfont
\setstretch{1.3}
\justifying

% ================================================================
% I. REPORT STRUCTURE
% ================================================================

\noindent
\textbf{I.\quad REPORT STRUCTURE}

\vspace{0.15cm}

The Final Design report shall be prepared using the official
\LaTeX{} template provided by the Department. Students shall not modify
the document structure, page layout, heading styles, numbering
system, or other predefined formatting settings unless approved by
the supervisor.

The report shall be arranged in the following order:

\begin{enumerate}[label=\arabic*., leftmargin=1.2cm, itemsep=2pt]
    \item Outer cover page;
    \item Inner cover page;
    \item Final Design Assignment Form;
    \item Supervisor's Evaluation Form;
    \item Acknowledgements (if applicable);
    \item Table of Contents;
    \item List of Symbols and Abbreviations (if applicable);
    \item List of Figures;
    \item List of Tables;
    \item Abstract;
    \item Main chapters of the FInal Design;
    \item Conclusions and Recommendations;
    \item References;
    \item Appendices (if applicable).
\end{enumerate}

The preliminary pages shall not be numbered as chapters. The main
content shall be organized into consecutively numbered chapters,
beginning with Chapter 1.

\vspace{0.3cm}

% ================================================================
% II. CONTENT OF THE FInal Design
% ================================================================

\noindent
\textbf{II.\quad CONTENT OF THE FInal Design}

\vspace{0.15cm}

\begin{itemize}[leftmargin=0.8cm, itemsep=3pt]

    \item The report shall contain only information directly related
    to the assigned project topic.

    \item The content shall demonstrate the student's own work,
    understanding, analysis, design, calculation, simulation,
    implementation, and evaluation as appropriate to the project.

    \item Information obtained from books, journal papers,
    conference papers, standards, websites, technical documents,
    or other sources must be properly cited.

    \item Chapters, sections, and subsections shall be generated
    automatically using the standard \LaTeX{} commands:

\begin{verbatim}
\chapter{...}
\section{...}
\subsection{...}
\end{verbatim}

    \item Chapter, section, and subsection numbers shall not be
    manually entered.

    \item The recommended hierarchy is Chapter -- Section --
    Subsection. Excessive subdivision of the report should be avoided.

    \item Unless otherwise specified by the Faculty or supervisor,
    the main report should normally contain between 30 and 80 pages,
    excluding large appendices.

\end{itemize}

\vspace{0.3cm}

% ================================================================
% III. PAGE LAYOUT AND TEXT FORMATTING
% ================================================================

\noindent
\textbf{III.\quad PAGE LAYOUT AND TEXT FORMATTING}

\vspace{0.15cm}

The report shall be formatted on A4 paper and printed on one side
of each sheet. The page margins shall be:

\begin{center}
\begin{tabular}{ll}
    Left margin:   & 3.5 cm \\
    Right margin:  & 2.0 cm \\
    Top margin:    & 2.5 cm \\
    Bottom margin: & 2.5 cm
\end{tabular}
\end{center}

The main text shall use a font size of 13 pt with 1.3-line spacing.
A Times New Roman or equivalent Times-family typeface shall be used.

The text shall be fully justified. Students shall not use repeated
spaces, manual line breaks, or repeated blank lines to position
text.

Page layout and general formatting shall be controlled by the
official template. Students shall not manually alter margins, font
sizes, line spacing, heading spacing, or paragraph formatting for
individual pages.

\vspace{0.3cm}

% ================================================================
% IV. CHAPTERS, SECTIONS, AND SUBSECTIONS
% ================================================================

\noindent
\textbf{IV.\quad CHAPTERS, SECTIONS, AND SUBSECTIONS}

\vspace{0.15cm}

\begin{itemize}[leftmargin=0.8cm, itemsep=3pt]

    \item Chapter titles shall be generated using:

\begin{verbatim}
\chapter{Chapter Title}
\end{verbatim}

    \item Section titles shall be generated using:

\begin{verbatim}
\section{Section Title}
\end{verbatim}

    \item Subsection titles shall be generated using:

\begin{verbatim}
\subsection{Subsection Title}
\end{verbatim}

    \item Numbers such as Chapter 1, 1.1, 1.1.1, etc.\ shall be
    generated automatically by \LaTeX{}.

    \item Students shall not manually reproduce heading styles using
    commands such as \verb|\textbf|, \verb|\fontsize|,
    or \verb|\centering| instead of the appropriate structural
    command.

\end{itemize}

\vspace{0.3cm}

% ================================================================
% V. PAGE NUMBERING
% ================================================================

\noindent
\textbf{V.\quad PAGE NUMBERING}

\vspace{0.15cm}

Cover pages and official forms at the beginning of the report shall
not display page numbers.

The main body of the report shall use Arabic page numbers beginning
from Chapter 1.

Page numbers shall be generated automatically by \LaTeX{}.
Students shall not manually enter page numbers.

\vspace{0.3cm}

% ================================================================
% VI. TABLE OF CONTENTS AND LISTS
% ================================================================

\noindent
\textbf{VI.\quad TABLE OF CONTENTS AND LISTS}

\vspace{0.15cm}

The Table of Contents shall be generated automatically using:

\begin{verbatim}
\tableofcontents
\end{verbatim}

The List of Figures shall be generated using:

\begin{verbatim}
\listoffigures
\end{verbatim}

The List of Tables shall be generated using:

\begin{verbatim}
\listoftables
\end{verbatim}

Students shall not manually create or edit page numbers in the
Table of Contents, List of Figures, or List of Tables.

The document shall be compiled sufficiently to ensure that all page
numbers, references, and lists are correctly updated.

\vspace{0.3cm}

% ================================================================
% VII. FIGURES, DIAGRAMS, AND GRAPHS
% ================================================================

\noindent
\textbf{VII.\quad FIGURES, DIAGRAMS, AND GRAPHS}

\vspace{0.15cm}

\begin{itemize}[leftmargin=0.8cm, itemsep=3pt]

    \item Figures shall be clear, legible, and directly relevant to
    the report content.

    \item Figures shall be inserted using the standard
    \verb|figure| environment.

    \item Figure captions shall be placed below the figure.

    \item Figures shall be numbered automatically according to the
    numbering system defined in the template.

    \item Each figure that is referenced in the text shall be assigned
    a unique label.

\end{itemize}

Example:

\begin{verbatim}
\begin{figure}[H]
    \centering
    \includegraphics[width=0.8\textwidth]
                    {figures/system_model.pdf}
    \caption{System model}
    \label{fig:system-model}
\end{figure}
\end{verbatim}

Figures shall be referenced in the text using:

\begin{verbatim}
Figure~\ref{fig:system-model}
\end{verbatim}

Students shall not manually type figure numbers.

Figures reproduced or adapted from other sources shall include an
appropriate citation.

Whenever possible, vector graphics or high-resolution images should
be used. Low-quality screenshots should be avoided.

\vspace{0.3cm}

% ================================================================
% VIII. TABLES
% ================================================================

\noindent
\textbf{VIII.\quad TABLES}

\vspace{0.15cm}

Tables shall be prepared using appropriate \LaTeX{} environments
such as \verb|table|, \verb|tabular|, \verb|tabularx|, or
\verb|longtable|.

Table captions shall be placed above the table.

Example:

\begin{verbatim}
\begin{table}[H]
    \centering
    \caption{System parameters}
    \label{tab:system-parameters}

    \begin{tabular}{ccc}
        \hline
        Parameter & Symbol & Value \\
        \hline
        Resistance & $R$ & 0.7 $\Omega$ \\
        Inductance & $L$ & 2.72 mH \\
        \hline
    \end{tabular}
\end{table}
\end{verbatim}

Tables shall be referenced using:

\begin{verbatim}
Table~\ref{tab:system-parameters}
\end{verbatim}

Students shall not manually enter table numbers.

Tables should be created directly in \LaTeX{} rather than inserted
as screenshots whenever practicable.

\vspace{0.3cm}

% ================================================================
% IX. EQUATIONS AND MATHEMATICAL EXPRESSIONS
% ================================================================

\noindent
\textbf{IX.\quad EQUATIONS AND MATHEMATICAL EXPRESSIONS}

\vspace{0.15cm}

Mathematical expressions shall be typeset directly using \LaTeX{}.
Images or screenshots of equations shall not be used.

Short mathematical expressions appearing within a sentence shall
use inline mathematical mode.

Example:

\begin{verbatim}
$R = 0.7\,\Omega$
\end{verbatim}

Important equations shall be placed in a numbered equation
environment.

Example:

\begin{verbatim}
\begin{equation}
    J\frac{d\omega}{dt}
    = T_e - T_L - B\omega
    \label{eq:mechanical-model}
\end{equation}
\end{verbatim}

Equations shall be referenced using:

\begin{verbatim}
Equation~\eqref{eq:mechanical-model}
\end{verbatim}

Equation numbers shall not be manually entered.

Mathematical notation shall be consistent throughout the entire
report.

\vspace{0.3cm}

% ================================================================
% X. SYMBOLS, UNITS, AND ABBREVIATIONS
% ================================================================

\noindent
\textbf{X.\quad SYMBOLS, UNITS, AND ABBREVIATIONS}

\vspace{0.15cm}

\begin{itemize}[leftmargin=0.8cm, itemsep=3pt]

    \item SI units shall be used wherever applicable.

    \item Symbols and units shall be written consistently throughout
    the report.

    \item The \verb|siunitx| package should be used for physical
    quantities and units where appropriate.

\end{itemize}

Examples:

\begin{verbatim}
\SI{220}{\volt}
\SI{2.72}{\milli\henry}
\SI{450}{\radian\per\second}
\end{verbatim}

Symbols and abbreviations frequently used in the report shall be
included in the List of Symbols and Abbreviations.

Each abbreviation should normally be defined when it first appears
in the text.

\vspace{0.3cm}

% ================================================================
% XI. SOURCE CODE
% ================================================================

\noindent
\textbf{XI.\quad SOURCE CODE}

\vspace{0.15cm}

Source code that is necessary for explaining an algorithm,
controller, simulation, or implementation may be included in the
main report.

The \verb|listings| environment, or another environment defined by
the official template, shall be used.

Example:

\begin{verbatim}
\begin{lstlisting}[language=Matlab]
A = [-R/L 0;
      0  -R/L];

K = hinfsyn(P,nmeas,ncont);
\end{lstlisting}
\end{verbatim}

Long source-code listings that are not essential to the immediate
discussion should be placed in the appendices.

Source code should not normally be inserted as screenshots.

\vspace{0.3cm}

% ================================================================
% XII. CROSS-REFERENCING
% ================================================================

\noindent
\textbf{XII.\quad CROSS-REFERENCING}

\vspace{0.15cm}

Chapters, sections, figures, tables, equations, and other numbered
elements shall be cross-referenced using the standard \LaTeX{}
mechanism.

The following commands shall be used where appropriate:

\begin{verbatim}
\label{...}
\ref{...}
\eqref{...}
\end{verbatim}

For example:

\begin{verbatim}
As shown in Figure~\ref{fig:response}, ...
\end{verbatim}

Students shall not manually type a reference number such as
``Figure 3.2'' when the number can be generated automatically.

Cross-referencing ensures that all numbers remain correct when
figures, tables, equations, or sections are added or removed.

\vspace{0.3cm}

% ================================================================
% XIII. REFERENCES AND CITATIONS
% ================================================================

\noindent
\textbf{XIII.\quad REFERENCES AND CITATIONS}

\vspace{0.15cm}

All sources used in preparing the FInal Design shall be properly
cited.

References shall be managed using the bibliography system provided
with the official template. Bibliographic information shall normally
be stored in the file:

\begin{verbatim}
References.bib
\end{verbatim}

A reference may be cited in the text using:

\begin{verbatim}
\cite{reference-key}
\end{verbatim}

Students shall not manually enter reference numbers such as
``[1]'', ``[2]'', etc.

Bibliographic entries shall contain sufficient information to
identify the original source.

Titles of books, journal articles, conference papers, standards,
and other publications in foreign languages shall be retained in
their original language and shall not be translated solely for the
reference list.

All ideas, data, figures, tables, equations, algorithms, or results
obtained from another source shall be appropriately acknowledged.

\vspace{0.3cm}

% ================================================================
% XIV. APPENDICES
% ================================================================

\noindent
\textbf{XIV.\quad APPENDICES}

\vspace{0.15cm}

Supporting material that is useful but too extensive for the main
body of the report should be placed in the appendices.

Appendices may include:

\begin{itemize}[leftmargin=0.8cm, itemsep=2pt]
    \item complete source code;
    \item detailed calculation sheets;
    \item large data tables;
    \item technical drawings;
    \item additional simulation results;
    \item equipment specifications;
    \item supplementary design documents.
\end{itemize}

Appendices shall be created using the appendix environment defined
in the official template.

Example:

\begin{verbatim}
\begin{appendices}

\chapter{MATLAB SOURCE CODE}

...

\chapter{SYSTEM PARAMETERS}

...

\end{appendices}
\end{verbatim}

Appendix titles and numbering shall not be manually entered.

\vspace{0.3cm}

% ================================================================
% XV. LATEX PROJECT ORGANIZATION
% ================================================================

\noindent
\textbf{XV.\quad \LaTeX{} PROJECT ORGANIZATION}

\vspace{0.15cm}

Students shall use the official \LaTeX{} template provided by the
Faculty.

The main formatting commands contained in the document preamble
shall not be modified unless approval is obtained from the
supervisor.

A recommended project structure is:

\begin{verbatim}
FinalDesign/
|
|-- main.tex
|-- References.bib
|
|-- figures/
|   |-- figure01.pdf
|   |-- figure02.png
|   `-- ...
|
|
`-- appendix/
    `-- ...
\end{verbatim}

File and folder names should preferably contain Latin letters,
numbers, hyphens, or underscores.

Spaces, Vietnamese diacritics, and unnecessary special characters
in file names should be avoided where possible.

Recommended examples include:

\begin{verbatim}
motor_model.pdf
chapter_01.tex
simulation_result.png
\end{verbatim}

\vspace{0.3cm}

% ================================================================
% XVI. DOCUMENT COMPILATION AND FINAL CHECK
% ================================================================

\noindent
\textbf{XVI.\quad DOCUMENT COMPILATION AND FINAL CHECK}

\vspace{0.15cm}

Before submission, the complete \LaTeX{} project shall compile
successfully and produce the final PDF report.

The final document shall not contain unresolved references such as:

\begin{verbatim}
??
[?]
\end{verbatim}

Before generating the final PDF, students shall check:

\begin{itemize}[leftmargin=0.8cm, itemsep=2pt]
    \item the Table of Contents;
    \item page numbering;
    \item chapter and section numbering;
    \item figure numbering;
    \item table numbering;
    \item equation numbering;
    \item cross-references;
    \item citations;
    \item bibliography;
    \item appendices;
    \item font display;
    \item mathematical symbols;
    \item image quality.
\end{itemize}

The document may be compiled using TeX Live, MiKTeX, or
another compatible \LaTeX{} environment. Regardless of the
compilation environment, the resulting PDF shall conform to the
official Department template. \textbf{Students are not allowed to use Overleaf 
or other online \LaTeX{} editors to protect the 
knowledge and intellectual property of the project.}

\vspace{0.3cm}

% ================================================================
% XVII. SUBMISSION REQUIREMENTS
% ================================================================

\noindent
\textbf{XVII.\quad SUBMISSION REQUIREMENTS}

\vspace{0.15cm}

The printed report shall satisfy the following requirements:

\begin{itemize}[leftmargin=0.8cm, itemsep=2pt]
    \item A4 paper;
    \item single-sided printing;
    \item one bound copy, unless otherwise specified;
    \item soft cover;
    \item formatting in accordance with the official template.
\end{itemize}

The electronic submission should include:

\begin{itemize}[leftmargin=0.8cm, itemsep=2pt]
    \item the final PDF report;
    \item the complete \LaTeX{} source files;
    \item the \verb|References.bib| file;
    \item all figures and graphics;
    \item all additional files required to compile the document.
\end{itemize}

The submitted source project shall be complete and capable of being
recompiled to produce a PDF corresponding to the submitted report.

\vspace{0.3cm}

% ================================================================
% XVIII. ELIGIBILITY FOR FInal Design DEFENSE
% ================================================================

\noindent
\textbf{XVIII.\quad ELIGIBILITY FOR FINAL DESIGN DEFENSE}

\vspace{0.15cm}

A student shall be eligible to defend the FInal Design when the
following requirements have been satisfied:

\begin{enumerate}[label=\arabic*., leftmargin=1.2cm, itemsep=2pt]

    \item The FInal Design has been completed within the specified
    period.

    \item The report satisfies the academic and technical requirements
    of the assigned project.

    \item The report complies with the official \LaTeX{} template
    issued by the Faculty.

    \item The complete \LaTeX{} source compiles successfully without
    errors affecting the document content or presentation.

    \item The final PDF and required source files have been completed
    and submitted.

    \item All external sources have been appropriately cited and the
    reference list has been properly generated.

    \item The supervisor has approved the student to proceed to the
    FInal Design defense.

\end{enumerate}

\vspace{0.4cm}

\noindent
\textbf{Note:}
Students are responsible for the academic and technical content of
the report. The official \LaTeX{} template is responsible for the
standardized document formatting. Students should therefore focus
on entering and organizing the content correctly rather than
manually modifying the predefined formatting parameters. 

\noindent
\textbf{Note:}
This section is provided for guidance only and is included in the main
document using the command \verb|\input{Regulations.tex}|.
Students shall remove or comment out this command before submitting the
final report and shall prepare the report content in accordance with the
requirements specified in this template.


} % End fontsize

\clearpage|.
Students shall remove or comment out this command before submitting the
final report and shall prepare the report content in accordance with the
requirements specified in this template.


} % End fontsize

\clearpage|.
Students shall remove or comment out this command before submitting the
final report and shall prepare the report content in accordance with the
requirements specified in this template.


} % End fontsize

\clearpage|.
Students shall remove or comment out this command before submitting the
final report and shall prepare the report content in accordance with the
requirements specified in this template.


} % End fontsize

\clearpage